\documentclass[11pt]{article}
\usepackage[letterpaper,margin=1in]{geometry}
\usepackage[T1]{fontenc}
\usepackage[hidelinks]{hyperref}
\pagestyle{empty}
\setlength{\parindent}{0pt}
\setlength{\parskip}{0.6em}
\begin{document}
\hfill 2 November 2026

Editor-in-Chief\\IEEE Transactions on Communications

Dear Editor-in-Chief,

I submit for your consideration ``MMSE Pair Scheduling: Computational Complexity and Certified Power Allocation,'' coauthored with Pooya Hejazi.

We study uplink scheduling with a common linear-MMSE receiver, two users per orthogonal resource, service of every user once, and continuous powers subject to individual bounds and an aggregate budget. The question is how tightly the quality of a feasible schedule can be bounded when pairing and powers change the receiver.

The paper establishes restricted hardness with three receive dimensions and a fixed rational channel alphabet, a sharp pair-rate concavity condition, and a characterization of rectangular pricing by at most seventeen candidates. It also develops independently checkable continuous upper bounds. Network pricing uses verified rectangular subdivision; the finite-candidate oracle is tested separately.

In two ray-traced scenes, pricing tightens 11 of 24 eight-user bounds and all eight sixteen-user bounds. It adds no achieved-rate gain over the corresponding search without pricing, and no main or scale case reaches the registered 0.01-nat target. The paper reports these limits and a channel-window error with its separately specified correction.

We cite and distinguish our earlier studies of clustering and base-station association with SIC, and of LS initialization with GAN-augmented CNN channel estimation. The present contribution concerns no-SIC scheduling and bounds for fixed nominal channels. Its focus on resource allocation, receiver-dependent optimization and performance limits fits the journal's communications scope.

The package includes editable sources, figures, supplementary proofs and results, and code with saved-evidence checking instructions. Raw DeepMIMO scenarios are linked separately, and AI assistance is disclosed in the manuscript. I will serve as corresponding author.

Thank you for considering our manuscript.

Sincerely,\\
Masoud Sabaei\\
Corresponding author\\
Department of Computer Engineering\\
Amirkabir University of Technology, Tehran, Iran\\
\href{mailto:sabaei@aut.ac.ir}{sabaei@aut.ac.ir}
\end{document}
